\begin{abstract}
We prove that every finite nonempty graph $G$ satisfies
$\chi_{\mathrm{list}}(G)\le C h(G)$ for an absolute integer $C$, where
$h(G)$ is the largest order of a clique minor. This resolves the Linear
List Hadwiger conjecture affirmatively.
\end{abstract}
\section{Introduction}\label{sec:introduction}

A clique minor records a form of complete interaction inside a graph:
its vertices are represented by disjoint connected sets, and every two
sets have an edge between them. For a finite nonempty simple graph $G$,
let $h(G)$ be the largest order of a clique minor in $G$. Hadwiger
conjectured that the chromatic number always satisfies
$\chi(G)\le h(G)$ \cite{Hadwiger1943}. The companion manuscript
\cite[Corollary~1.2]{OpenAIHadwigerCounterexample2026} constructs finite
simple graphs with $\chi(G)>h(G)$, disproving this ordinary coefficient-one
assertion. Its linear relaxation asks whether
an absolute constant can replace the coefficient one. We prove a bound
of this form even when allowable colors vary from vertex to vertex.

A \emph{list assignment} $\mathcal L$ assigns a finite set of colors to
each vertex. An $\mathcal L$-coloring is a proper coloring $c$ satisfying
$c(v)\in\mathcal L(v)$ at every vertex. The list chromatic number
$\ell(G)=\chi_{\mathrm{list}}(G)$ is the least integer $k$ for which every
assignment with $|\mathcal L(v)|\ge k$ is colorable. List coloring was
introduced independently by Vizing \cite{Vizing1976} and by Erd\H{o}s,
Rubin, and Taylor \cite{ErdosRubinTaylor1980}. The linear list
version of Hadwiger's conjecture was formulated by Kawarabayashi and
Mohar \cite[Conjecture~1.3]{KawarabayashiMohar2007}.

\begin{theorem}\label{thm:main}
There is an absolute integer $C\ge1$ such that every finite nonempty
simple graph $G$ satisfies
\[
 \ell(G)\le C h(G).
\]
\end{theorem}

Thus the Linear List Hadwiger conjecture has an affirmative resolution.
The constant is independent of the graph order, the list assignment, and
the number of distinct colors appearing in its union. Taking the same
list at every vertex immediately gives the ordinary conclusion.

\begin{corollary}\label{cor:ordinary}
Every finite nonempty simple graph $G$ satisfies $\chi(G)\le C h(G)$,
with the constant $C$ in Theorem~\ref{thm:main}.
\end{corollary}
\begin{proof}
A coloring from the common list $\{1,\ldots,\ell(G)\}$ shows that
$\chi(G)\le\ell(G)$.
\end{proof}
